% Transcription — fonds Alexandre Grothendieck, shelfmark 29, batch 11 (pages 201-216).
% Established from the facsimile archives/batches/29/batch-11.pdf (a mirror of
% https://grothendieck.umontpellier.fr/29.pdf), per the skill
% transcribe-grothendieck.
%
% Pass: Opus 5 (claude-opus-5), 2026-08-29 — first pass, unchecked against the
% pages by a human.
%
% This is the last batch of the shelfmark: sixteen pages, of which fifteen are
% transcribed. Page 207 is a blank leaf and gets no \page{}.
%
% The batch falls into three runs.
%
%   Pages 201-206 — the end of the "domaines de ramification" run the wrapper of
%     page 188 opened, which batch 10 left mid-sentence at page 200. He
%     paginates it 11 to 16, continuing his 1-10 of batch 10. The subject is
%     generating families: a fibred category C over Et(S) with a functor
%     φ: C → R, the subcategories R_S^{C_0} attached to the admissible sieves
%     C_0 of C, the functor α into the inverse limit, and the proof that α is
%     fully faithful. The run stops at page 206 on "Prouvons que α est
%     essentiellement surjectif", and the proof does not follow.
%
%   Pages 208-212, 214-216 — a fresh run, "Catégories multigaloisiennes",
%     written out fair in blue ink and much the most legible thing in the
%     batch. He paginates it 1, 2, 3, 4 (pages 208-211), then 1 bis (page 212)
%     and 1 ter (page 214) — two inserted leaves carrying the "4 conditions sur
%     feuille 1 bis" that page 208 promises under A(iv) — then 5 and 6 (pages
%     215, 216). His numbers on pages 212 and 214-216 are circled.
%
%   Page 213 — a landscape sheet on quite another subject: the fibred powers
%     X^E of a morphism f: X → Y, the partitions P of E, the closed immersions
%     α_P, and a count of dimensions in the case card E = 3. It carries no
%     number of his, and the whole sheet is crossed by two long diagonals. It
%     is transcribed where it stands, in page order, and no continuity with
%     what surrounds it is asserted.
%
% Notation fixed once for the batch, following batch 10: Et(S) for the étale
% site, Fl_S for the fibred category of schemes over it. Three readings had to
% be settled and are flagged where they first occur: the subscript "cf" of
% T_cf and C̃_cf, which he glosses himself as "constants finis"; the script E
% of (E|C_0), which he never names; and the circled S of the symmetric group
% on page 209. The four headings of page 212 are numbered backwards, 3, 2, 1,
% 0, each digit circled — rendered here as (3), (2), (1), (0).
%
% The dating is the inventory's own, for the whole shelfmark.
\documentclass[11pt,a4paper]{article}
% Le préambule commun à toutes les transcriptions du fonds Grothendieck.
%
% Il tient en une idée : **l'incertitude se compose, elle ne se résout pas.**
% Une transcription qui lisse un mot douteux a détruit la seule chose pour
% laquelle on la faisait — un lecteur doit pouvoir savoir, à chaque endroit,
% ce qui était sur le feuillet et ce qui a été ajouté. Les six macros qui
% suivent sont donc l'appareil critique, et non de la décoration.
%
% Les mêmes commandes sont reconnues par `scripts/render.mjs`, qui produit la
% vue de lecture du site. Une macro ajoutée ici doit l'être aussi là-bas,
% faute de quoi le rendu échouera bruyamment — ce qui est voulu : un rendu
% silencieusement faux est pire qu'un rendu absent.

\ProvidesPackage{grothendieck}[2026/08/08 Transcriptions du fonds Grothendieck]

\usepackage{amsmath,amssymb,amsthm}
\usepackage{mathrsfs}
\usepackage{stmaryrd}
% Les diagrammes commutatifs sont partout dans ce fonds : c'est la langue
% même de la géométrie algébrique de Grothendieck, et une transcription qui
% les rendrait en prose ne transcrirait rien.
\usepackage{tikz-cd}
% Français par défaut — c'est la langue des feuillets — mais l'anglais est
% chargé aussi : la traduction et le résumé sont des documents anglais, et la
% typographie française (espaces avant les deux-points) y serait une faute.
% Ces documents basculent par \selectlanguage{english} en tête de corps.
\usepackage[english,french]{babel}
\usepackage{fontspec}
\usepackage{geometry}
\geometry{a4paper,margin=25mm,marginparwidth=28mm}
\usepackage{marginnote}
\usepackage{xcolor}
% ulem, not soul. `soul`'s \st and \ul are built on a character-by-character
% scan that XeLaTeX's font machinery breaks: the document fails with an
% undefined control sequence rather than degrading. ulem does the same two jobs
% and is Unicode-engine safe. `normalem` stops it hijacking \emph, which the
% transcriptions use for Grothendieck's own emphasis.
\usepackage[normalem]{ulem}
\usepackage{hyperref}
\hypersetup{colorlinks=true,linkcolor=black,urlcolor=black,pdfborder={0 0 0}}

% --- Métadonnées du lot -----------------------------------------------------
% Reprises telles quelles par le script de rendu, qui les affiche en tête de
% la vue de lecture. Elles fixent ce que le fichier prétend couvrir : sans
% elles, une transcription flotte sans point d'attache dans un fonds de seize
% mille feuillets.

\newcommand{\folder}[1]{\gdef\grothFolder{#1}}
\newcommand{\batch}[1]{\gdef\grothBatch{#1}}
\newcommand{\leaves}[2]{\gdef\grothFirst{#1}\gdef\grothLast{#2}}
\newcommand{\foldertitle}[1]{\gdef\grothFoldertitle{#1}}
\gdef\grothFolder{?}\gdef\grothBatch{?}\gdef\grothFirst{?}\gdef\grothLast{?}\gdef\grothFoldertitle{}

% --- Le résumé --------------------------------------------------------------
%
% Il ouvre la lecture modernisée, sous le titre « Résumé », et il est composé
% en petit. Ce n'est pas un ornement : c'est le seul passage du document qui
% ne soit pas des mathématiques, et un lecteur doit pouvoir le reconnaître
% avant de l'avoir lu — puis le sauter s'il connaît déjà le sujet, ou s'y
% arrêter s'il ne le connaît pas. Le filet qui le suit marque la reprise du
% corps.

\newenvironment{resume}
  {\small\setlength{\parskip}{0.45em}}
  {\par\medskip\hrule\medskip}

% --- Mots-clés ---------------------------------------------------------------
%
% En anglais, en clôture du résumé de la lecture modernisée : le vocabulaire
% moderne sous lequel chercher ce que les feuillets construisent. Le manifeste
% du site (`npm run manifest`) les extrait pour étiqueter la cote entière —
% cette ligne est la source unique des tags, il n'y a pas de fichier à côté.
\newcommand{\keywords}[1]{\par\smallskip\noindent
  {\footnotesize\textsc{Keywords} --- \textit{#1}}\par}

% --- L'appareil critique ----------------------------------------------------

% Le numéro de feuillet, en marge. C'est l'ancre qui permet de régler un
% désaccord sur une formule en regardant *un* feuillet plutôt qu'une chemise
% de six cents. Le site s'en sert aussi pour tourner les pages du fac-similé
% quand on fait défiler la transcription.
\newcommand{\leaf}[1]{%
  \marginnote{\footnotesize\color{gray!70}\textsf{#1}}%
  \label{leaf:#1}\ignorespaces}

% Illisible : on ne devine pas. Un mot inventé qui se lit comme les autres est
% le pire résultat possible.
\newcommand{\ill}{\textcolor{red!60!black}{[\,\ldots\,]}}

% Lecture douteuse : le mot est donné, et signalé comme incertain.
\newcommand{\uncertain}[1]{\uline{#1}}

% Ajout de l'éditeur — une lettre manquante, un mot sous-entendu.
\newcommand{\add}[1]{\textcolor{blue!50!black}{[#1]}}

% Biffé par Grothendieck lui-même. Ce qu'il a rayé fait partie du document :
% une rature dit souvent où la pensée a changé de direction.
\newcommand{\struck}[1]{\sout{#1}}

% Note du transcripteur, sur l'état matériel du feuillet.
\newcommand{\note}[1]{\par\smallskip{\small\color{gray!60!black}\textit{[#1]}}\par\smallskip}

% Note marginale de Grothendieck, distincte du corps du texte.
\newcommand{\marginal}[1]{\par\smallskip\noindent\hspace{1em}%
  {\small\color{gray!60!black}#1}\par\smallskip}

% --- Titre ------------------------------------------------------------------

\AtBeginDocument{%
  \begin{center}
    {\large\bfseries Fonds Alexandre Grothendieck --- cote n\textsuperscript{o}~\grothFolder}\\[2pt]
    {\itshape\grothFoldertitle}\\[4pt]
    {\small feuillets \grothFirst--\grothLast\ (lot \grothBatch)}\\[2pt]
    {\footnotesize\color{gray!60!black}Transcription établie d'après le fac-similé numérisé
     par l'Université de Montpellier.\\
     Les lectures incertaines sont soulignées, les illisibles notées [\,\ldots\,],
     les ajouts entre crochets.}
  \end{center}
  \vspace{1em}}

\endinput


\folder{29}
\batch{11}
\pages{201}{216}
\dating{1959-1969}
\watermark{Édition de démonstration}
\foldertitle{Groupe fondamental [Autour de SGA 4 et SGA 7] : lettres (1959, 1967, 1969), tapuscrits et copies de tapuscrit annotés (s.d.), notes manuscrites (s.d.)}

\begin{document}

\section*{Familles génératrices d'un domaine de ramification (pages 201 à 206)}

\note{la suite ouverte à la page 190 se poursuit sans rupture ; il la pagine
lui-même 11 à 16, en haut à droite, à la suite du 1-10 du lot précédent. La
page 200 s'arrêtait en cours de phrase, sur « Dans le cas de l'exemple 2 » ; ce
qui suit n'y revient pas et repart sur les familles génératrices}

\page{201}\note{sa page 11}

Soit $(R, r)$ un domaine de ramification sur $S$. Une famille d'objets
$X_i/S_i$ de $R$ ($S_i \in \mathrm{ob}\,\mathrm{Et}(S)$) est dite
\emph{génératrice} si pour tt objet $X$ de $R$ sur un $S'$ de
$\mathrm{Et}(S)$, les $S''$ étales sur $S'$ tels qu'il existe $i$ et
$S'' \to S_i$ et \struck{\ill{}} $X_{S''}$ \emph{subordonné} à $X_{i,S''}$,
recouvrent $S'$. Si \struck{les} \add{de plus} \uncertain{adaptés} \ill{}
$X_i$ \ill{} générateur \ill{} $R_{S_i}$, \struck{\ill{}} que $G_i$, d'où
$(Z_i, G_i)$ \ill{} : générateurs \ill{} $r(R_i)$ sur $S_i$, les $Z_i$, $G_i$
doivent permettre de reconstituer $(R, r)$, une fois précisé les \ill{}
\ill{} de vérifier, mais la formulation de telles données est pénible et
finalement peu utile.

\note{la phrase se défait entre « Si de plus » et « les $Z_i$, $G_i$ doivent
permettre » ; ce qui se lit est donné tel quel, et ce qui ne se lit pas est
marqué}

Il est plus commode de \struck{\ill{}} associer à $R$ une \struck{catégorie
fibrée} \add{catégorie} $C$ sur $\mathrm{Et}(S)$, et un foncteur
$\varphi : C \longrightarrow R$. On dit que $\varphi$ est \emph{génératrice}
si les $\varphi(M)$ ($M \in \mathrm{ob}\,C$) forment une \emph{famille}
génératrice, i.e. si pour tt $X \in \mathrm{ob}\,R$ sur un
$S' \in \mathrm{ob}\,\mathrm{Et}(S)$, les $S''$ étales sur $S'$ tels que
$X_{S''}$ soit \emph{subordonné} à un objet de $C_{S''}$
\add{soit $\varphi(M) \to X_{S''}$ \ill{}}, \struck{forment un} recouvrement
de $S'$. Posant $s = r\varphi$
\[
  s : C \longrightarrow \mathrm{Fl}_S
\]
\ill{} dit qu'on peut alors, à équivalence près,

\page{202}\note{sa page 12}

reconstituer $(R, r)$ en termes de $(C, s)$. Bien sûr on peut également
reconstituer $(R_{S'}, r_{S'})$, la construction s'appliquer donc aux
$R_{S'}$, $r_{S'}$. (En s'induisant sur $S'$ et applique ce qu'on peut avoir
\ill{} fait à $R|S'$) et les fonctorialités seront aussi évidentes\,\ldots

Soit $X$ un objet de $R_S$. Soit $C_X$ la ss.-catégorie pleine de $C$ formée
des objets $M$ \add{i.e. pour $S' \in \mathrm{ob}\,\mathrm{Et}(S)$}
\struck{\ill{}} tels que $X_{S'}$ soit subordonné à $\varphi(M)$ \add{et tels
que de plus $\varphi(M) \to X_{S'}$ soit un épim.}, \ill{} $M$
\struck{\ill{}} \add{admissible de $C$} \ill{} par hypothèse ; et \ill{}

Alors pour tout \struck{tel} $M \in \mathrm{ob}\,C_X$,
$r(X_{S'} \times \varphi(M)) \longrightarrow r\varphi(M) = s(M)$ est un
\struck{morphisme} \add{revêtement} étale. \struck{Donc} Ainsi
\[
  M \longmapsto r(X \times_{S'} \varphi(M)) \big/ \underbrace{r(\varphi(M))}_{s(M)}
\]
donne une \emph{section cartésienne} de la catégorie fibrée sur $C_X$
\struck{induit par $s : C \to$} \add{sur $\mathrm{Fl}_S$} (catégorie fibrée
\struck{dont} \add{les} revêtements étales relatifs) via
$C_X \xrightarrow{\ s\ } \mathrm{Fl}_S$.

\note{le foncteur composé $C_X \to \mathrm{Fl}_S$ est écrit avec le $s$
au-dessus de la flèche}

\struck{\ill{} la catégorie fibrée} obtenue \add{ainsi} \struck{sur} $C$ via
$s : C \to \mathrm{Fl}_S$, \struck{\ill{}} \ill{} quel \uncertain{crible} de
$C_0$, \struck{\ill{} $R_S^{C_0}$} \struck{fait} la sous-catégorie
\struck{fibrée} \add{pleine} de $R_S$ formée des objets $X$
\struck{\ill{}} tels que $\forall\, S' \in \mathrm{ob}\,\mathrm{Et}(S)$,
\struck{\ill{}} $X_{S'}$ est \add{loc.} \emph{subordonné} à tt objet de $C_0$
sur $S'$. \ill{} trouve donc une ss.-catégorie \uncertain{multigaloisienne} de
$R_S$, stable \ill{} les $\varprojlim$ et sommes directes, et un foncteur
\[
  (\ast) \qquad \alpha : R_S^{C_0} \longrightarrow
  \varprojlim_{C_0} \, (\mathcal{E}|C_0)
\]

\note{le symbole restreint aux cribles est un E de fantaisie ; il ne le nomme
nulle part, et il est rendu ici par $\mathcal{E}$ partout où il paraît}

\page{203}\note{sa page 13}

Pour $C_0$ parcourant les cribles \struck{couvrants} \add{admissibles} de $C$,
\struck{les $R$} et admettant que la catégorie \add{fibrée} $C$ est
\emph{localement filtrante} \ill{}, ce qui implique que l'intersection de deux
cribles couvrants est encore couvrant, on voit que les $R_S^{C_0}$ forment une
famille filtrante de ss.-catégories multigaloisiennes pleines de $R_S$, de
\ill{} réunion $R_S$. Donc \ill{} si on voit que les foncteurs $\alpha_0$ sont
des \emph{équivalences de catégories}, on aura alors une équivalence
\[
  R_S \;\simeq\; \varinjlim_{C_0} \; \varprojlim_{\text{sur } C_0} \,
  (\mathcal{E}|C_0)
\]
qui explicite $R_S$ en termes de $(C, s)$, \add{mais} \struck{mais} qui reste
\struck{au moins} que la notion de \emph{crible admissible}. On va montrer
\ill{} pour que les \struck{\ill{}} \ill{} il suffit que $\varphi$ soit
\emph{pleinement fidèle}.

Comme $\alpha$ est \ill{} un foncteur exact de catégories multigaloisiennes,
il suffit de voir pour le \struck{plus} précis que
\[
  \alpha(X) = \emptyset \;\Longrightarrow\; X = \emptyset ,
\]
ce qui est immédiat \ill{} $C_0$ \add{admissible} \struck{couvrant} : c'est
local par $(\mathrm{Et})$, \struck{aux $\varphi(M) \to e_S$ épim. \ill{}}
donc \ill{} il existe un $M \in C_{0S}$, \ill{}
$r(\varphi(M) \times X) = \emptyset \Longrightarrow \varphi(M) \times X =
\emptyset$. Ok.

$\varphi$ pleinement fidèle ? Soit $X \in \mathrm{ob}\,R_S^{C_0}$, et
$Z = \alpha(X) = Z_1 \sqcup Z_2$, montrons que cette

\marginal{\ill{} ; \ill{} $\longrightarrow C$ ; loc. filtrante}

\note{cette dernière note est cerclée, en bas à gauche de la page}

\page{204}\note{sa page 14}

décomposition est induite par une décomposition de $X$. La question étant
locale sur $S$, \ill{} peut supposer $\exists\, M \in \mathrm{Ob}\,C_S$,
$\varphi(M) \longrightarrow e$ un épimorphisme.

Considérons $r(X \times \varphi(M)) = \underbrace{\alpha(X)(M)}_{= Z(M)}
\longrightarrow r(\varphi(M)) = s(M)$ ; il se coupe en
$Z_1(M) \sqcup Z_2(M)$, donc \struck{par} \textbf{R2}
$X \times \varphi(M)$ se coupe en $X_1 \sqcup X_2$. Je dis que
\struck{\ill{}} cette décomposition provient d'une décomposition de $X$, i.e.
que quels que soient les $C$-objets \struck{$T$} \add{$T$} de $R_S^{C_0}$ et
deux morphismes $T \rightrightarrows \varphi(M)$, les deux images inverses de
$X_1$ par $T \rightrightarrows \varphi(M)$ sont les \ill{}, i.e.
\struck{\ill{}}

\begin{tikzcd}[column sep=small, row sep=small]
T \arrow[d, bend left=12] \arrow[d, bend right=12]
  & X \times T \arrow[d, bend left=12] \arrow[d, bend right=12] & \\
\varphi(M) \arrow[d]
  & X \times \varphi(M) \arrow[l, no head] \arrow[r] \arrow[d] & X_1 \\
e & X \arrow[l, no head] &
\end{tikzcd}

donc $T \times \varphi(M) \longrightarrow \varphi(M)$ un épim. Comme
$T \times \varphi(M)$ est \emph{subordonné} à $\varphi(M)$,
$T \times \varphi(M)$ est \emph{constant par morceaux} sur $\varphi(M)$, et
\ill{} $=$ épimorphique, il est \emph{recouvrant} (pour sa topologie
canonique, v. cor.) par les images de sections \add{$w$} de
$T \times \varphi(M)$ sur $\varphi(M)$. Ceci permis, on peut remplacer
$T \rightrightarrows \varphi(M)$ par les composés
\[
  T \times \varphi(M) \xrightarrow{\ \mathrm{pr}_1\ } T
  \rightrightarrows \varphi(M)
\]
(ou plus un épim., $\varphi(M) \longrightarrow e_S$ \add{étant}\,!), puis en
dernier par les composés
$\varphi(M) \xrightarrow{\ w\ } T \times \varphi(M)$. Bref, on est ramené au
cas où $T = \varphi(M)$. Mais $\varphi$ étant pl. fidèle,

\marginal{on a $r(X_1) \rightrightarrows$ \ill{} $= r(X \times T)$ \ill{}}

\page{205}\note{sa page 15}

on trouve que les deux flèches sont identiques par
$M \rightrightarrows M$, et \struck{on gagne} \struck{\ill{}} de voir que les
deux images de $r(X_1)$ dans
$r(X \times T) = r(X \times_S \varphi(M)) = \alpha(X)(M)$ est $Z_1(M)$.

On a donc terminé une décomposition
\add{uniquement déterminée par la condition $Z_1 = \alpha(X_1)$}
\[
  X = X_1 \sqcup X_2 ,
\]
\struck{Je dis que} la décomposition correspondante de $\alpha(X)(M)$ est
donnée par $Z_1(M)$.

\note{à partir d'ici et jusqu'au bas de la page, le texte est accolé à gauche
et barré de deux longues diagonales ; il se laisse lire par endroits et par
endroits non}

i.e. \ill{} que \ill{} tt $M'$, et si $M' \to$ \ill{}, $Z'_1 = \alpha(X_1)$,
\ill{} $Z_1$ et $Z'_1$, \ill{} $A$ et $B$. \ill{}
$Z_1 \cap Z'_1$ dans $Z_1$ et $Z'_1$, soient $A$ et $B$. \ill{}
\[
  A(M) = \emptyset , \quad B(M) = \emptyset ,
  \qquad A(M') = \emptyset , \quad B(M') = \emptyset
\]
\ill{} pour $M' \in C_{0S}$, \ill{}
$A(M) = \emptyset \Longrightarrow M = \emptyset$. \ill{} en effet,
\struck{\ill{}} $M'$ \struck{\ill{}} $A(M') \neq \emptyset$ (car il \ill{} : à
vérifier :) \ill{} on peut supposer $S' = S$, et \ill{}
$A(M \times M') = \emptyset$, donc \ill{}
\[
  A(M \times M') = A(M') \times_{s(M')} s(M \times M') ,
\]
et \add{[on voit que $s(M \times M') \longrightarrow s(M')$ \struck{est
surjectif} ou encore \ill{}]} $(M \times) M' \longrightarrow M$, c'est
trivial. Finalement, il \ill{} localiser, il existe $M''$, $M'' \to M$,
$M'' \to M'$.

\begin{tikzcd}[column sep=small, row sep=small]
& \varphi(M) \times_{\varphi} M' \arrow[dl] \arrow[dr] & \\
\varphi(M) & & \varphi(M')
\end{tikzcd}

Cas. $A(M'') = \emptyset$, $B(M'') = \emptyset$. Donc la décomposition
\ill{} aussi celles qui seront définies \ill{} lieu des \ill{} $M$. Mais il y
a \ill{} $M'$, et \ill{} $X = X'_1 \sqcup X'_2$, \ill{} celle

\page{206}\note{sa page 16}

un autre $M'$, et si $\nexists\, M' \longrightarrow M$, alors il est \ill{}
que la décomposition de \struck{$X$} définie par $M'$ et la \ill{}.
Localisant sur $S$ et utilisant le fait que $C$ est loc. filtrante, on trouve
que la décomposition \emph{ne dépend pas du choix de $M$}. Localisant à
nouveau sur $S$, on voit qu'elle \ill{} la décomposition donnée de $Z$.

Prouvons que $\alpha$ est \emph{essentiellement surjectif}.

\note{la page s'arrête là ; la démonstration annoncée ne suit pas, et la suite
de la cote est d'un autre sujet}

\section*{Catégories multigaloisiennes : les conditions A et B, les axiomes G0 à G5 (pages 208 à 212)}

\note{la page 207 est nue. Ce qui suit est écrit au net, à l'encre bleue, et
se lit presque de bout en bout, à la différence de tout ce qui précède. Il
pagine la suite lui-même : 1 à 4 pour les pages 208 à 211, puis 1 bis et 1 ter
--- deux feuillets insérés --- pour les pages 212 et 214, puis 5 et 6 pour les
pages 215 et 216. Les numéros des pages 212 et 214 à 216 sont cerclés}

\page{208}\note{sa page 1}

\textbf{Catégories multigaloisiennes.}

$\mathcal{U}$ un univers, $C$ une catégorie.

\textbf{A) Conditions équivalentes}

\begin{enumerate}
\item[(i)] Il existe un $\mathcal{U}$-topos $\mathcal{T}$, et une équivalence
$C \simeq \mathcal{T}_{\mathrm{cf}}$, où $\mathcal{T}_{\mathrm{cf}}$ est la
catégorie pleine de $\mathcal{T}$ formée des objets loc.\textsuperscript{t}
« constants finis ».

\item[(ii)] \add{$C$ est une $\mathcal{U}$-catégorie, l'ens. des objets
\struck{\ill{}} mod. iso. de $C$ est $\mathcal{U}$-petit, et} [ Il existe une
topologie sur $C$, [moins fine que la canonique] et telle que le foncteur de
$C \longrightarrow \widetilde{C}$ induise une équivalence
$C \longrightarrow \widetilde{C}_{\mathrm{cf}}$.

\item[(ii bis)] Id., en prenant la « top. la plus fine de la descente
effective universelle ».

\item[(iii)] [ $C$ est une $\mathcal{U}$-catégorie, dont l'ens. des objets
mod. iso. est $\mathcal{U}$-petit ] Pour tout objet $X$ de $C$,
l'ens\struck{emble} \struck{des} objets $S$ tels que $X_S$ sur $S$ soit
\emph{constant fini} sur $S$ forment une famille \struck{couvrante} de
descente effective universelle.

\item[(iv)] 4 conditions sur feuille 1 bis.
\end{enumerate}

\note{l'indice des deux catégories est le même et se lit « cf » ; il le
glose lui-même à la ligne au-dessus, « constants finis », et c'est cette
glose qui l'a fixé ici. Les deux crochets ouvrants de (ii) et (iii) sont sur
la page et ne se referment pas}

\textbf{B) Conditions équivalentes}

\note{le bloc qui suit est accolé à gauche et barré de deux diagonales ; il
est remplacé par le (i) qui le suit}

\begin{enumerate}
\item[(ii)] \struck{Comme (ii) dessus, mais en exigeant de plus que l'objet
final de $\mathcal{T}$ est somme finie d'objets connexes.}

\item[] \textbf{(ii bis), (iii), (iv)} \struck{Comme plus haut, mais en exigeant que
l'objet final de $C$ (ou encore, \ill{} \add{d'un \ill{}}) est somme finie
d'objets connexes.}
\end{enumerate}

\begin{enumerate}
\item[(i)] $C$ satisfaisant aux conditions de A), et son objet final est somme
\emph{effective universelle} d'objets connexes.
\end{enumerate}

\page{209}\note{sa page 2}

\begin{enumerate}
\item[(ii)] On a les conditions \textbf{G0} à \textbf{G5} :
\end{enumerate}

\begin{itemize}
\item[\textbf{G0}] $C$ est une $\mathcal{U}$-catégorie, l'ens. des
\struck{classes} \struck{\ill{}} \emph{classes d'isomorphie} d'objets de $C$
est $\mathcal{U}$-petit.

\item[\textbf{G1}] Existence de $\varprojlim$ \add{inv.} \emph{finies} et
\struck{$\varinjlim$} \add{sommes finies}.

\item[\textbf{G2}] Les \struck{épim.} \emph{relations d'équivalence} sont
effectives, \struck{les épim. \ill{} effectifs universels}
\add{\struck{finies sont disjointes}} $\mid$ épi. universelles.

\item[\textbf{G3}] Tout morphisme $X \longrightarrow Y$ se factorise en un
\struck{morphisme de descente effective} \struck{épimorphisme}, suivi d'un
isom. sur un \uncertain{sommand} direct.

\item[\textbf{G4}] \struck{Tout} L'objet \add{final} \struck{est} \struck{une}
somme \struck{disjointe} \emph{effective univ.} d'objets connexes.

\item[\textbf{G5}] \struck{\ill{} [ou encore : l'objet final est réunion finie
d'objets connexes]} Soit $f : X \longrightarrow Y$ avec $X$, $Y$ connexes ;
alors $\exists$ entier $N$ tel que pour tt $Y' \longrightarrow Y$,
$Y' \neq \emptyset_C$, \struck{$X \times_Y Y'$ \ill{} connexes} \add{\ill{}}
\end{itemize}

\begin{enumerate}
\item[(iii)] \struck{\ill{}} $C$ est équivalente à une catégorie produit
$\prod_{i \in I} C_i$, les $C_i$ \emph{galoisiennes}.
\end{enumerate}

[ \textbf{N.B.} les $C_i$ correspondent \ill{} une composante connexe
\add{$e_i$} de l'objet final ; la donnée d'une équivalence
$C_i \simeq C(\pi_i)$ correspond à la donnée d'un \emph{foncteur fibre}
\add{de $C_{/e_i}$} d'un « rev. universel » de l'objet $e_i$. ]

\note{un signe précède ce dernier $e_i$ que je ne lis pas}

\textbf{Démonstration de A}

(iii) $\Rightarrow$ (ii bis) $\Rightarrow$ (ii) $\Rightarrow$ (i) trivial ; et
(i) $\Rightarrow$ (iii) aussi ; car si $X \in \mathcal{T}_{\mathrm{cf}}$,
\struck{\ill{}} $\mathcal{T}_{\mathrm{cf}}$, \struck{est} l'objet final de
$\mathcal{T}$ est somme (donc dans $\mathcal{T}$) d'objets $e_i$
($i \in \mathbb{N}$) aux $X|e_i$ de rang $i$. De plus,
$\mathbf{Isom}_{e_i}\bigl([1, i]_{e_i} ,\, X\bigr) = P_i \longrightarrow$
\struck{dans} un objet de $\mathcal{T}_{\mathrm{cf}}$ tel que $P_i \longrightarrow e_i$ épim. (i.e.
un \emph{torseur} sous $(\mathfrak{S}_i)_{e_i}$\,!) et tel que

\note{le groupe est écrit d'un S cerclé ; l'énoncé --- un torseur sous lequel
un objet de rang $i$ devient constant --- le donne pour le groupe symétrique,
et il est rendu ici $\mathfrak{S}_i$}

\page{210}\note{sa page 3}

$X_{P_i} \simeq [1, i]_{P_i}$. Les $P_i \longrightarrow e$ forment une famille
\struck{épimorphique} de descente effective universelle pour les
$\mathcal{T}$, donc aussi \struck{des} $\mathcal{T}_{\mathrm{cf}}$
( \textbf{N.B.} $\mathcal{T}_{\mathrm{cf}}$ \struck{stable} \add{est} A.B.
par produits finis \struck{et} \struck{même} \add{et par $\varprojlim$ et
$\varinjlim$ finis} ).

\textbf{Démonstration de B} \quad En fait, on constate aisément que les
conditions de A(ii) impliquent \textbf{G0} \textbf{G1} \textbf{G2}
\textbf{G3}.

\note{ce qui suit jusqu'au bas de la page est barré d'une longue diagonale ;
seuls des fragments s'en laissent lire}

\struck{L'équivalence de (i) et (iii) tient au fait que les décompositions
\struck{en somme} de $e$, dans $\mathcal{T}_{\mathrm{cf}}$ \struck{ou} dans
$\mathcal{T}$, \struck{sont les \ill{}} donc disent que \ill{} dans $C$,
\ill{}}

\struck{Pour prouver que les conditions impliquent (iv), il \struck{faut}
\add{prouver} \ill{} qu'elles impliquent \textbf{G4}, i.e. tout revêtement
\struck{fini} d'un objet connexe $e$ est somme finie d'objets connexes. En
fait, si $X \longrightarrow e$ de rang $n$ sur $e$, \ill{} le \ill{} de
composantes connexes de $X$ est $\leq n$, \struck{et ce qui} les
$X \longrightarrow$ \struck{$e$} finis : c'est immédiat par récurrence sur $n$
et « définition du rang ». \ill{} --- prenant pour $\mathcal{T}$ le topos
$(B\pi_i)$ \ill{}}

\struck{(v) $\Rightarrow$ (iii) \ill{} « somme » des $B\pi_i$ où chaque
$\pi$ \ill{} \ill{} prouver (iv) $\Rightarrow$ (iii). Je préfère (ii)
$\Rightarrow$ (iii).}

\page{211}\note{sa page 4}

Pour \ill{} \struck{exiger} \textbf{G0} : \textbf{G4}, il faut prouver ceci :
si $X \in \mathrm{ob}\,C$, alors il \struck{existe} \add{les}
\struck{une famille de descente effective universelle} $S_i \longrightarrow e$
telles que les $X_{S_i}/S_i$ soient \emph{constant finis}, recouvrant $e$ par
la topologie de la descente effective. Utilisant \textbf{G2} et \textbf{G4},
\ill{} \add{sommant} sur $e_i$ ( somme \struck{non} disjointe « (c) » et
produit fibrés constants ), \ill{}

\marginal{[\textbf{G3}]}

\note{ce qui suit est accolé à gauche et barré de deux longues diagonales}

\struck{\ill{} et universelle, et forment une famille sur $e'$ le \ill{}
« \ill{} certain ». Je dis, En effet, considérons \ill{} $e'$ \ill{} $e''$ son
« complémentaire » \ill{} $e'' = \emptyset$ i.e. $e' = e$.}

\struck{$X_{e''}$ sur $e''$, alors} \ill{} que si $X \longrightarrow Y$ ($Y$
\struck{non} \ill{} non vide \ill{}) est tel que $Y$ \ill{} $X$ connexe, alors
$\exists\, Y' \longrightarrow Y$ \ill{} universelle, avec
$X_{Y'} \simeq [1, n]_{Y'}$ : procéder par récurrence sur l'entier
\struck{en baissant les dimensions} (\ill{} le rang], on \ill{} constatation
plus haut \ill{} \textbf{G5} \ill{}

\textbf{Remarque} : il se peut \ill{} Notons qu'en prisant \ill{}
\textbf{G5}$'$ \struck{\ill{}} \ill{} $\exists$ foncteur exact \ill{}
prouvant que \textbf{G5} est \emph{redondant} \ill{} des autres axiomes, on
\ill{} voir : pour tt composant connexe $e_i$ de $e$, \ill{}
\[
  F : C \longrightarrow (\mathrm{Ens}\ f) \quad \text{tel que}
  \quad F(e_i) \neq \emptyset
\]

\page{212}\note{sa page 1 bis, cerclée ; c'est la « feuille 1 bis » que la
page 208 annonce sous A(iv). Ses quatre rubriques sont numérotées à rebours,
3, 2, 1, 0, chaque chiffre cerclé}

$\exists$ \emph{foncteur}

\textbf{(3)} On les définit, pour tt \ill{} $A$ \ill{} $\neq \emptyset$, la
classe des « morphismes de degré $n$ », de façon que

\begin{enumerate}
\item[a)] \struck{stabilité} par changement \ill{}
\item[b)] \emph{additivité} \emph{finie}
\item[c)] Un isom. \ill{} est de degré $1$
\item[\struck{d)}] Si $f : X \longrightarrow Y$, \ill{} alors
$Y = \displaystyle\coprod_{n \in \mathbb{N}} Y_n$, somme directe disjointe,
effective, universelle, tels que
$X_n = X \times_Y Y_n \longrightarrow Y_n$ soit de degré $n$.
\item[d)] Si $f : X \longrightarrow Y$ est de degré $n \geq 1$, alors $f$
\ill{} un morphisme de descente \emph{effective}
[ donc \emph{de descente effective universelle} grâce à a) ]
\item[e)] Si $f : X \longrightarrow Y$ est de degré $0$, alors
$X = \emptyset_C$ \ill{} et un isomorphisme avec \ill{} \emph{sommand direct
universel}.
\end{enumerate}

\textbf{(2)} Tous les monomorphismes \ill{} \ill{} un \ill{}

\textbf{(1)} Existence des $\varprojlim$ \emph{finies} \ill{}

\struck{\textbf{(\ill{})} $e$ \ill{}}

\textbf{(0)} $C$ est une \emph{$\mathcal{U}$-catégorie} (i.e. les
$\mathrm{Hom}(X, Y)$ sont $\mathcal{U}$-petits) et \struck{l'on} \struck{des}
\ill{} classes d'isom. d'objets de $C$ \ill{} $\mathcal{U}$-petites.

$'$ Indiquer condition A(iii). En effet, prouvons que pour $X$ \ill{} donné,
les $S \to e$ tels que $X_S$ constant fini sur $S$ forment \ill{} famille de
descente eff. universelle. Ainsi : a), f) \ill{} on est ramené à prouver
l'analogue pour $X \longrightarrow Y$ \ill{} $\varprojlim$ \ill{} avec
$X \longrightarrow Y$ de degré $n$ \struck{( à trouver un morphisme de
descente eff. universelle qui trivialise )}. Si $n = 0$, on conclut par 3 a)
puis $X = \emptyset$, \ill{} prend $Y \longrightarrow Y$
($X = \emptyset_Y$\,!). Si $n > 0$, on procède par récurrence supposant que

\section*{Puissances fibrées, partitions et comptes de dimension (page 213)}

\note{feuillet en travers, d'un tout autre sujet que ce qui l'entoure : les
puissances fibrées $X^E$ d'un morphisme $f : X \to Y$, les partitions $P$ de
$E$, et un compte de dimensions. Il ne le pagine pas, et le feuillet entier
est barré de deux longues diagonales. Il est transcrit ici à son rang, sans
qu'aucune continuité avec la suite multigaloisienne ne soit affirmée}

\page{213}

\begin{tikzcd}[column sep=large]
X^P \arrow[r, hook, "\alpha_P"] & X^E
\end{tikzcd}

\note{sous $\alpha_P$ il écrit « imm. fermée » ; sous $X^P$ un trait crochu
porte $F_P$, sans tête, et son sens n'est pas fixé}

\[
  X^E = \bigcup_P \alpha_P\bigl((X^P)^{\ast}\bigr)
  \qquad \text{(disjoint)}
\]
\[
  (f^E)^{-1}\bigl(\Delta_y^{(E)}\bigr) = \bigcup_P \alpha_P(F_P)
  \qquad \text{(réunion disjointe)}
\]
de dim $n - (\mathrm{card}\,P - 1)d \;\geq\; n - \mathrm{card}(E)\,d$, égalité
pour $P = P_0$ seulement $\hookrightarrow$ \emph{partition maximale}.

$\underline{\mathrm{card}\,E = 3}$ :

\begin{itemize}
\item partition en 3 : $P_0$
\item partitions en 2 : $(1)(23)$, $(2)(31)$, $(3)(12)$ \qquad
$P_{23}$, $P_{31}$, $P_{12}$
\item partition en 3 : $P_1$
\end{itemize}

\note{les deux lignes extrêmes portent le même chiffre, 3, là où l'une des
deux devrait porter 1 ; le chiffre est écrit deux fois de la même main et
laissé tel quel --- \emph{sic}}

\[
  (f^3)^{-1}\bigl(\Delta_y^{(3)}\bigr) =
  \underset{n - 2d}{F_3} \;\cup\;
  \underset{n - d}{\bigcup_{1 \leq i < j \leq 3} \alpha_{ij}(F_2)} \;\cup\;
  \underset{n}{\Delta_X^{(3)}}
\]

\marginal{dim « génériques »}

\note{suit un bloc barré, sous le mot « Question » lui-même barré ; il se lit
en partie}

\struck{$\overline{F}_3 \cap \Delta_X^{(3)} = \emptyset$ ;
$\overline{\alpha_{ij}(F_2)} \cap \Delta_X^{(3)} =
\alpha_{ij}(\overline{F}_2) \cap \Delta_X^{(3)} =
\alpha_{ij}\bigl(\overline{F}_2 \cap \Delta_X^{(3)}\bigr) = \emptyset$ ;
$\overline{F}_3 \cap \alpha_{ij}(F_2)$,
$\overline{F}_3 \cap \alpha_{ij}(\overline{F}_2)$ \ill{}}

\textbf{N.B.} Si $f$ est \ill{}, alors
$\overline{F}_3 \cap \Delta_X^{(3)} = \emptyset$,
$\overline{\alpha_{ij}(F_2)} \cap \Delta_X^{(3)} = \emptyset$.

\marginal{fermé}

Mais que peut-on dire de
\[
  \overline{F}_3 \cap \overline{\alpha_{ij}(F_2)} \;\cdot\;
  \overline{F}_3 \cap \alpha_{ij}(\overline{F}_2) \;\;?
\]

\[
  x_1^0 \neq x_2^0 \neq x_3^0 \qquad\qquad x_1^0 = x_3^0 \neq x_2^0
\]

\note{les exposants de ces six points sont un seul trait rond ; $0$ est la
lecture retenue, sans certitude}

\begin{tikzcd}[column sep=small, row sep=small]
X(1) \arrow[dr] & & X(2) \arrow[dl] \\
& y &
\end{tikzcd}

\[
  \dim X = 3 , \quad \dim Y = 4 , \qquad \dim F_3 = 1 , \qquad
  \dim \alpha_{ij}(F_2) = 2
\]

\note{en bas à gauche, un petit croquis --- deux ouverts $U$ et $V$ se
coupant, avec « $x_1, x_2 \in U$ » --- qui n'est pas repris ici}

\section*{Morphismes de degré $n$, les conditions C, et la proposition finale (pages 214 à 216)}

\page{214}\note{sa page 1 ter, cerclée ; le feuillet fait suite à la page 1
bis}

c'est prouvé pour les entiers $< n$. \{ Utilisant 3 d),

\begin{tikzcd}[column sep=small, row sep=small]
X \arrow[d] & X \times_Y X = X' \arrow[l] \arrow[d] \\
Y & X = Y' \arrow[l, no head]
\end{tikzcd}

et regarder l'\struck{épimorphisme} \emph{diagonal} $\Delta_{X,Y}$ ( existe
par \textbf{(1)} ), par \textbf{(2)} \ill{} un \emph{sommand direct}
\add{universel} donc $X' \simeq Y' \sqcup X'_1$ ; par 3 a)
$X' \longrightarrow Y'$ \ill{} est de degré \struck{$n$} \add{$n-1$}, par 3,
donc par 3 b) et f) \struck{$X'$} $\longrightarrow Y'$ \ill{} de degré 1, on
voit que $X'_1 \longrightarrow Y'$ est de degré $n-1$. Les \ill{} \ill{}
l'hyp. de récurrence, on trouve $Y''_s \longrightarrow Y'$ famille
\struck{\ill{}} de descente eff. universelle, qui trivialise $X'_1$ ; comme
$X' \simeq Y' \sqcup X'_1$ \emph{universellement},
\[
  X'' = X \times_Y Y'' \simeq Y'' \sqcup
  \underbrace{(X'_1 \times_{Y'} Y'')}_{Y'' \times [1,\, n-1]}
  = Y'' \times [1, n] .
\]
$Y'' \longrightarrow Y$ est de descente \ill{} effective universelle, comme
composé de deux tels, OK.

\textbf{Explicitation.} Pour vérifier \ill{} qu'un $f$ \ill{} morphisme
\ill{} et de descente effective [cf. condition 3 d)], il suffit de savoir que
c'est un \emph{épimorphisme effectif universel}, et de savoir que les
\emph{relations d'équivalence sont effectives universelles} \ill{}

\begin{tikzcd}[column sep=small, row sep=small]
Z'' \arrow[d]
  & Z \arrow[l, dashed] \arrow[d]
  & D \arrow[l, bend left=12] \arrow[l, bend right=12] \arrow[d] \\
Y & X \arrow[l]
  & X \times_Y X \arrow[l, bend left=12] \arrow[l, bend right=12]
\end{tikzcd}

\note{le nom du sommet supérieur gauche porte deux traits surmontés de points
et n'est pas assuré}

\textbf{Conclusion} \quad Dans les conditions de A(iv), la notion de
« morphisme de degré $n$ » est \emph{uniquement déterminée}, et
\struck{correspond} équivaut à : $\exists\, Y' \longrightarrow Y$
épimorphisme \struck{\ill{}} effectif universel, et de descente effective
\ill{}, avec $X \times_Y Y' \simeq [1, n]_{Y'}$ ( $Y' \to$ \ill{} isom.
\ill{} ).

\page{215}\note{sa page 5, cerclée}

\textbf{C. Conditions équivalentes}, \ill{} \emph{galoisiennes}

\begin{enumerate}
\item[(i)] Comme A(i), avec l'objet final de $\mathcal{T}$ connexe.

\item[] \textbf{(ii), (ii bis), (iii), (iv)} Comme dans A, en supposant de plus l'objet
final \ill{} de $C$ connexe.

\item[(v)] Comme A(iv) : les axiomes \textbf{G0} \textbf{G1} \textbf{G2}
\textbf{G3}, et de plus : $\exists$ foncteur \struck{conservatif}
\[
  F : C \longrightarrow (\mathrm{Ens\ fini})
\]
(qui est \emph{conservatif}), \emph{exact} : gardant et commutant aux sommes
finies.

\item[(vi)] \struck{\ill{}} Il existe un \emph{groupe profini} $\pi$ et une
équivalence $C \simeq C(\pi)$.
\end{enumerate}

\page{216}\note{sa page 6, cerclée}

\textbf{Proposition} \quad Soit $u : C \longrightarrow C'$ un foncteur
\emph{exact}, avec $C$, $C'$ des catégories multigaloisiennes. Alors :

\begin{enumerate}
\item[(i)] $u$ \emph{fidèle} $\Longleftrightarrow$ $u$ \emph{conservatif}
$\Longleftrightarrow$ pour tt objet $X$ de $C$,
$u(X) = \emptyset \Longrightarrow X = \emptyset$.

\item[(ii)] $u$ \emph{pleinement fidèle} $\Longleftrightarrow$ pour tt objet
$X$ de $C$, \struck{toute} l'application \struck{\ill{}} \ill{}
\struck{$X$} $u(X)$ est une \emph{bijection} de l'ens. des sous-objets de $X$
\struck{avec} l'ens. des ss.-objets de $Y$.

\item[(iii)] $u$ \emph{une équivalence} $\Longleftrightarrow$ $u$ est pl.
fidèle (cf. (ii)) et tout objet de $C'$ est \emph{subordonné} à un objet de
\ill{} $u(X)$ \ldots
\end{enumerate}

\marginal{si $u$ n'est \struck{\ill{}} pl. fidèle, \struck{\ill{}}, tt \ill{}
ss.-objets et \ill{} par les \ill{} subordonnés}

\note{la moitié inférieure de la page est nue ; la cote s'arrête là}

\end{document}
